% =====================================================================
%  Chapter 33 — The Selection Tests
% =====================================================================

\chapter{The Selection Tests}
\label{ch:selection-results}

Chapter 22 built the four tests. This chapter applies them to the book's
own family and reports what they say. They say that the system's return
is not distinguishable from the benchmark's, and that the machinery of
selecting among the family is, on this evidence, worse than a coin flip.

\section{The family}

The trial family\idx{Trial family} has \numFamilySize{} members: \numModelConfigs{} model
configurations, three ablations that turn each layer off, two Kelly
fractions, two tilt scales, two volatility targets, two ensemble
half-lives, and two alternative strategic policies.

It is deliberately wide. A reality check\idx{Reality check} over twelve near-identical
model configurations would pass trivially and would mean nothing; a check
over a family that includes \emph{turn the machine off} and \emph{use a
completely different allocation} is a check that could fail.

\section{What the tests say}

\begin{center}
\begin{tabular}{@{}P{1.72in}P{0.72in}P{2.30in}@{}}
\toprule
Test & Value & Reading \\
\midrule
Deflated Sharpe Ratio\idx{Deflated Sharpe ratio}       & \numDSR{} & Not significant. The observed Sharpe is below the expected maximum for a family this size. \\
DSR with implied trials     & \numDSRImplied{} & Not significant. The same conclusion on the correlated-trial correction. \\
Probability of backtest overfitting & \numPBO{} & Above one half. The selection procedure is worse than a coin flip. \\
White reality check         & p = \numRCP{} & No member of the family beats the index by more than chance. \\
Hansen superior predictive ability & p = \numSPAP{} & The same conclusion, with the family reweighted by its own performance. \\
Newey--West alpha against the index & \numAlphaT{} (p = \numAlphaP{}) & The excess return is \numAlpha{} a month, and is indistinguishable from zero. \\
\bottomrule
\end{tabular}
\end{center}

\begin{takeaway}[The verdict]
None of the five tests finds evidence that the system beats the index, and
the probability of backtest overfitting is above one half, which means the
procedure of selecting among the family's members is actively unreliable
on this evidence. The book reports this and does not reinterpret it.
\end{takeaway}

\section{Reading the probability of backtest overfitting}

The PBO deserves more than a line, because it is the most informative of
the five and the least familiar.

The construction: split the record into eight contiguous sub-periods in
every possible way --- \numPBOSplits{} splits --- and for each split,
rank the \numPBOTrials{} family members by their Sharpe ratio on the first
half, then see where the first-half winner lands in the second half's
ranking. The statistic is the probability that it lands in the bottom
half.

The book's value is \numPBO{}, which is above one half. The probability
that the in-sample winner is out-of-sample in the bottom half is
\numPBOProb{}, which is below one half.

The reading is unambiguous: selecting the configuration that worked best
in one half of the record gives you a configuration that is, more likely
than not, among the worse half on the other half of the record. The
selection procedure is doing worse than chance.

\section{Why the family is effectively one trial}

The Deflated Sharpe Ratio's most striking input is the number of
effectively independent trials, which the book computes from the family's
own correlation matrix.

The average pairwise correlation between the family's return series is
\numAvgCorr{} --- the members produce nearly identical portfolios --- and
the implied number of independent trials\idx{Implied number of independent trials} is therefore
\numImpliedTrials{}, rather than the \numFamilySize{} members actually run.

This has two consequences, and they pull in opposite directions.

The first is that the Deflated Sharpe Ratio's penalty is smaller than a
naive count would give. With a family of twenty-five near-identical
members, the expected maximum Sharpe ratio under the null is
\numSRZero{} --- and the book's observed Sharpe ratio is not above it. The
Monte Carlo estimate of the same quantity, from twenty thousand simulated
families, is \numSRZeroMonteCarlo{}.

The second is the more important one, and it is a methodological point
rather than a numerical one. \textbf{If the members of your research
family all produce nearly the same portfolio, then you did not try
twenty-five strategies; you tried one.} The trial count is a property of
the strategies, not of the number of configurations that generate them.

The family's spread against the index is shown in \figref{fig:family}.

\section{What the alpha test says}

The book's system earns an excess return over the index of \numAlpha{} a
month --- about \numAlpha{} a month over \numAlphaMonths{} months --- with a
HAC standard error of \numAlphaSE{}, giving a t-statistic of
\numAlphaT{} and a p-value of \numAlphaP{}.

The t-statistic is not close to significant, and the book would not claim
it was even if it were. The relevant comparison is not with the critical
value of a two-sided test at five per cent. It is with the critical value
that the Deflated Sharpe Ratio implies, which is a much higher bar, and
which this result does not reach.

\section{What the selection tests do not say}

Three things the reader should not take from this chapter.

They do not say the system is worthless. They say the system's excess
return is not \emph{distinguishable from zero}, which is a statement about
what the evidence can distinguish, not about what is true. A real edge of
half a percentage point a year would look exactly like this in a
twenty-four-year record with a family of this size.

They do not say the equal-weighted portfolio is better. The equal-weighted
portfolio is a comparator, not a member of the family, and it is not
subject to the selection tests. It earns more compound return and takes a
deeper drawdown, and the reader must choose between those, which is a
judgement and not a measurement.

They do not say the architecture is wrong. The tests say that the
\emph{prediction layer} does not contribute. They say nothing about the
governor, whose contribution was positive in every one of the
\numWindowCount{} starting windows in Chapter 31, because the governor
makes no forecast and therefore generates no selection problem.

\begin{keytakeaways}
\item The Deflated Sharpe Ratio is \numDSR{}, the probability of backtest
overfitting is \numPBO{}, and both the reality check and the superior
predictive ability test return a p-value of one.
\item The family's members are so similar that the implied number of
independent trials is \numImpliedTrials{} rather than \numFamilySize{},
and that is a methodological point rather than a technical one.
\item The Newey--West alpha against the index is \numAlphaT{} with a
p-value of \numAlphaP{}.
\item The tests say the prediction layer contributes nothing
distinguishable from zero; they say nothing about the governor, which
makes no forecast.
\end{keytakeaways}
